\subsection{Transporte afín de los coeficientes de Laurent}
\label{ix-add:stieltjes}

El cambio positivo y adimensional $Q\mapsto\lambda Q$, con
$\lambda>0$, actúa sobre el lector posterior mediante
\[
 Z_\lambda(s)=\lambda^{-s}Z(s).
\]
La identidad se verifica primero en $\Re s>1$ y continúa
meromórficamente. El residuo de $Z_\lambda$ en uno es $\lambda^{-1}$.
Para comparar coeficientes con residuo unitario se define
\[
 \widehat Z_\lambda(s)=\lambda Z_\lambda(s)
 =e^{-\ell(s-1)}Z(s),\qquad \ell=\log\lambda,
\]
\[
 \widehat Z_\lambda(s)=\frac1{s-1}
 +\sum_{n\geq0}\frac{(-1)^n\gamma_n(\lambda)}{n!}(s-1)^n.
\]

\HMTresultado{Teorema}{Ley de transporte de Stieltjes}{ix-add:stieltjes-law}
Para todo $n\geq0$,
\[
 \gamma_n(\lambda)=
 \sum_{j=0}^{n}\binom nj\gamma_j\ell^{n-j}
 -\frac{\ell^{n+1}}{n+1}.
\]
Esta transformación es triangular afín. Dos cambios de escala se
componen multiplicando sus parámetros y la escala inversa proporciona
la transformación inversa.

\textbf{Demostración.} Póngase $z=s-1$ y multiplíquese la serie
completa de Laurent por $e^{-\ell z}$. El término polar contribuye al
coeficiente de $z^n$ con $(-\ell)^{n+1}/(n+1)!$; los términos regulares
contribuyen con
\[
 \sum_{j=0}^n\frac{(-1)^j\gamma_j}{j!}
                    \frac{(-\ell)^{n-j}}{(n-j)!}.
\]
Multiplicar la suma por $(-1)^nn!$ produce la fórmula. El producto
$e^{-\ell z}e^{-mz}=e^{-(\ell+m)z}$ prueba la ley de composición.
La sustitución $m=-\ell$ da la inversa. \hfill$\square$

En particular,
\[
 \gamma_0(\lambda)=\gamma_0-\ell,\qquad
 \gamma_1(\lambda)=\gamma_1+\gamma_0\ell-\frac{\ell^2}{2}.
\]
El término afín registra la contribución del polo y ya interviene en el
primer coeficiente. Para $\lambda=3$ se recupera la transformación de
la clase residual cero: sus coeficientes son los de
$Z_0(s)=3^{-s}Z(s)=\widehat Z_3(s)/3$ de la
sección~\ref{lectura:manuscrito-04}. Esta especialización conserva
el factor $1/3$ del residuo. La ley general permite transportar la
normalización sin recalcular ni seleccionar otro lector primario TPK.
